\documentclass[11pt]{article}
\usepackage[a4paper,margin=1in]{geometry}
\usepackage{amsmath,amssymb,mathtools,bm}
\usepackage{enumitem,booktabs,array}
\usepackage{tcolorbox}
\usepackage{hyperref}
\usepackage[protrusion=true,expansion=false]{microtype}
\hypersetup{colorlinks=true,linkcolor=blue,urlcolor=blue,citecolor=blue}
\setlist[itemize]{topsep=2pt,itemsep=2pt}
\setlist[enumerate]{topsep=2pt,itemsep=2pt}
\newtcolorbox{keybox}{colback=blue!3!white,colframe=blue!40!black,boxrule=0.4pt,arc=1mm}
\newtcolorbox{alertbox}{colback=red!3!white,colframe=red!40!black,boxrule=0.4pt,arc=1mm}
\newtcolorbox{answerbox}{colback=green!3!white,colframe=green!40!black,boxrule=0.4pt,arc=1mm}
\newtcolorbox{drillbox}{colback=yellow!5!white,colframe=orange!60!black,boxrule=0.4pt,arc=1mm}
\newcommand{\EV}{\mathbb{E}}
\newcommand{\Var}{\mathrm{Var}}
\newcommand{\Cov}{\mathrm{Cov}}
\newcommand{\blank}[1]{\underline{\hspace{#1}}}

\title{W10-04: Fahlenbrach, Rageth \& Stulz (2021, RFS) \\
\large ``How Valuable is Financial Flexibility When Revenue Stops? Evidence from the COVID-19 Crisis'' --- Active Recall Workbook}
\author{Banking \& Risk Management --- Exam Prep}
\date{\today}
\begin{document}\maketitle

\begin{keybox}
\textbf{Paper in one sentence.} In the March--April 2020 COVID revenue shock, firms with greater ex-ante \emph{financial flexibility} (high cash-to-assets, low leverage, high credit-line availability) outperformed by $\sim$6--10 pp in cumulative stock returns --- revealing the \emph{option value} of liquidity in a genuine revenue-stop scenario.

\textbf{Placement.} The cleanest evidence on the stock-market pricing of financial flexibility. Complements BKS (W10-03) on precautionary cash and KRS (W01-01) on credit-line liquidity insurance.
\end{keybox}

\section*{Round 1: Complete Walk-Through}

\subsection*{1.1 Setting}
Mid-February through April 2020: broad market collapse as the COVID-19 pandemic triggered shelter-in-place orders. Revenue for many sectors \emph{stopped}. Financial-flexibility measures ex-ante (end of 2019): cash/assets, leverage, undrawn credit-line capacity.

\subsection*{1.2 Data}
US firm returns from CRSP. Financial-flexibility covariates from Compustat and Capital IQ. Sample: non-financial, non-utility US firms.

\subsection*{1.3 Empirical spec}
Cross-sectional regression of cumulative stock returns during the COVID crash on ex-ante flexibility measures plus controls:
\[
R_{i,\text{COVID}}=\alpha+\beta_1\cdot\text{Cash/Assets}_i+\beta_2\cdot\text{Leverage}_i+\beta_3\cdot\text{CreditLineCap}_i+X_i'\gamma+\varepsilon_i.
\]

\subsection*{1.4 Headline results}
\begin{itemize}
\item High-cash firms outperformed low-cash peers by $\sim$7 pp cumulative returns during the initial crash.
\item Low-leverage firms outperformed by $\sim$5--10 pp.
\item Undrawn credit-line capacity partially substitutes for cash but is weaker.
\item Effects concentrated in sectors hardest hit by revenue stops (travel, hospitality, retail).
\item The effect is robust to industry, size, and prior-year performance controls.
\end{itemize}

\subsection*{1.5 Interpretation}
\begin{itemize}
\item Financial flexibility has a large ex-post \emph{option} value that is under-priced in normal times.
\item FSS (1993) / precautionary-savings theory: cash and unused debt capacity let firms maintain operations and investment when external finance is unavailable.
\item BKS (2009): the upward trend in cash hoarding is rational precaution, as COVID shows.
\end{itemize}

\subsection*{1.6 Caveats}
\begin{alertbox}
\begin{itemize}
\item Single-event; generalisation requires care.
\item Reverse causality: firms anticipating trouble may have built up cash.
\item Ex-ante flexibility correlates with other firm traits (size, industry) --- controls important.
\end{itemize}
\end{alertbox}

\section*{Round 2: Level 1 Fill-in}
\begin{enumerate}
\item Shock window: \blank{1cm} through \blank{1cm} 2020.
\item Main flexibility measures: \blank{1cm}/assets, \blank{1cm}, credit-line capacity.
\item High-cash vs.\ low-cash return differential: $\sim$\blank{1cm} pp.
\item Sectors most affected: \blank{1.5cm}, hospitality, retail.
\item Theoretical anchor: \blank{1cm} precautionary savings.
\end{enumerate}

\section*{Round 3: Level 2 Prompts}
\begin{enumerate}
\item Explain why financial flexibility is an ``option'' on continuing operations.
\item Write the cross-sectional regression and state the identifying concern.
\item Discuss whether COVID generalises to other revenue-stop shocks.
\item Connect FRS to BKS (2009) and KRS (2002).
\end{enumerate}

\section*{Round 4: Minimal Scaffolding}
From a blank page: (i) shock, (ii) flexibility measures, (iii) regression, (iv) magnitudes, (v) interpretation.

\section*{Round 5: Blank Page}
Essay: what COVID teaches us about corporate financial flexibility.

\vspace{6pt}
\begin{drillbox}
\textbf{Speed Drill.} (1) Shock window. (2) Three flexibility measures. (3) Cash vs.\ no-cash differential. (4) Most-affected sectors. (5) Theoretical anchor paper.
\end{drillbox}

\section*{Quick Reference \& Answer Key}
\begin{answerbox}
\begin{itemize}
\item \textbf{Shock.} COVID-19 revenue stop, Feb--Apr 2020.
\item \textbf{Measures.} Cash, leverage, credit-line capacity.
\item \textbf{Effect.} High-flexibility firms outperform by $\sim$6--10 pp.
\item \textbf{Theory.} Option value of liquidity; FSS precautionary savings.
\end{itemize}
\end{answerbox}

\begin{keybox}
\textbf{30-second exam pitch.} Fahlenbrach, Rageth \& Stulz (2021) study the COVID-19 revenue shock as a natural experiment on the value of financial flexibility. In February--April 2020, when many US firms saw revenue collapse, those entering the shock with more ex-ante financial flexibility --- high cash-to-assets, low leverage, and available credit-line capacity --- dramatically outperformed peers in stock returns by roughly 6--10 percentage points cumulative. The return premium is largest in sectors with the sharpest revenue stops (travel, hospitality, retail). Read with Bates--Kahle--Stulz (2009), which documented the secular rise in corporate cash as a precautionary response to rising idiosyncratic risk, the COVID episode shows the option value is real. Flexibility is under-priced in normal times and acutely valuable in tail-risk states, validating the Froot--Scharfstein--Stein precautionary-savings logic of corporate liquidity.
\end{keybox}

\end{document}
